\documentclass[12pt]{article}
\usepackage[utf8]{inputenc}
\usepackage[T1]{fontenc}
\usepackage[margin=1in]{geometry}
\usepackage{setspace}
\usepackage{xcolor}
\usepackage[colorlinks=true, urlcolor=blue!60!black, linkcolor=blue!60!black]{hyperref}
\setlength{\parindent}{0em}
\setlength{\parskip}{0.45em}
\renewcommand{\baselinestretch}{1.0}
\begin{document}
\begin{center}
{\Large \textit{Hernandez v.\ Texas}, 347 U.S.\ 475 (1954)\par}
\vspace{0.3em}
{\normalsize Case Brief --- One-Page Edition\par}
\end{center}

\textbf{Memory Jogger.} Equal protection reaches beyond a Black--white ``two-class theory'': no systematic exclusion of Mexican Americans from juries (\href{https://supreme.justia.com/cases/federal/us/347/475/}{Justia}).

\textbf{Facts.} All-Anglo juries in Jackson County, Texas, indicted and convicted Pete Hernandez of murder. Texas stipulated that no person with a Mexican or Latin American name had served on a jury there in 25 years, out of more than 6,000 jurors, though qualified Mexican Americans lived in the county. Local segregation extended to a courthouse restroom marked ``Colored Men'' and ``Hombres Aqui.''

\textbf{Procedural History.} The Texas Court of Criminal Appeals affirmed: Mexican Americans were white, and the Amendment recognized only two classes.

\textbf{Votes.} Reversed, 9--0, Chief Justice Warren writing.

\textbf{Holding.} Excluding eligible jurors because of ancestry or national origin violates equal protection. The challenger must prove a distinct local class and a prima facie case of systematic exclusion under \textit{Norris v. Alabama}; general denials of intent do not rebut it.

\textbf{Reasoning.} Groups needing protection change with community prejudice. Community attitudes---testimony, school segregation, a ``No Mexicans Served'' sign---proved a class apart, and Spanish surnames identified its members. Zero jurors in 25 years could not be chance, whatever the commissioners' professed good faith. There were no separate opinions.

\textbf{Examples.} \emph{Same-Side:} no Hmong-American jurors in 20 years in a segregated county. \emph{Opposite-Side:} a random jury that happens to lack Mexican Americans. \emph{Fence-Sitter:} 30\% of residents but 4\% of grand jurors.

\textbf{Critique.} Haney L\'{o}pez hailed its view of race as a social construct, but proving class status county by county burdened Latino litigants and left the binary intact.

\textbf{Key Quotation.} ``Throughout our history, differences in race and color have defined easily identifiable groups which have at times required the aid of the courts in securing equal treatment under the laws. But community prejudices are not static, and, from time to time, other differences from the community norm may define other groups which need the same protection.'' --- Chief Justice Warren, 347 U.S. at 478.

\end{document}
